%%
%% Elsevier CAS Double Column Template for Computational Biology and Chemistry
%%
\documentclass[a4paper,fleqn]{cas-dc}

\usepackage[authoryear,longnamesfirst]{natbib}
\usepackage{amsmath,amssymb}
\usepackage{booktabs}
\usepackage{multirow}
\usepackage{algorithm}
\usepackage{algorithmic}
\usepackage{hyperref}
\usepackage{cleveref}
\usepackage{subcaption}
\usepackage{color}
\usepackage{url}

%%% Author macros
\def\tsc#1{\csdef{#1}{\textsc{\lowercase{#1}}\xspace}}
\tsc{WGM}
\tsc{QE}
%%%

% Custom commands
\newcommand{\ie}{\textit{i.e.}}
\newcommand{\eg}{\textit{e.g.}}
\newcommand{\etal}{\textit{et al.}}
\newcommand{\R}{\mathbb{R}}

\begin{document}
\let\WriteBookmarks\relax
\def\floatpagepagefraction{1}
\def\textpagefraction{.001}

% Short title
\shorttitle{[Short Title]}

% Short author
\shortauthors{[First Author et al.]}

% Main title of the paper
\title[mode=title]{[Your Paper Title Here]}

% Title footnote mark
\tnotemark[1]

% Title footnote
\tnotetext[1]{This research did not receive any specific grant from funding agencies in the public, commercial, or not-for-profit sectors.}

% First author
\author[1]{First Author}[%
  orcid=0000-0000-0000-0000,
  style=chinese
]

% Corresponding author indication
\cormark[1]

% Footnote of the first author
\fnmark[1]

% Email id of the first author
\ead{first.author@university.edu}

% URL of the first author
\ead[url]{https://university.edu/~firstauthor}

% Credit authorship
\credit{Conceptualization, Methodology, Software, Writing - original draft}

% Address/affiliation
\affiliation[1]{%
  organization={University Name},
  addressline={Address},
  city={City},
  postcode={000000},
  state={State},
  country={Country}
}

% Second author
\author[2]{Second Author}

% Footnote of the second author
\fnmark[2]

% Email id of the second author
\ead{second.author@university.edu}

% Credit authorship
\credit{Data curation, Validation, Writing - review \& editing}

% Address/affiliation
\affiliation[2]{%
  organization={Another University},
  addressline={Address},
  city={City},
  postcode={000000},
  state={State},
  country={Country}
}

% Corresponding author text
\cortext[1]{Corresponding author.}

% Footnote text
\fntext[1]{}

% Abstract
\begin{abstract}
[Your abstract here -- approximately 200--250 words. Describe the problem, proposed method, experimental setup, and key results.]
\end{abstract}

% Research highlights
\begin{highlights}
\item [Highlight 1: key contribution or finding]
\item [Highlight 2: methodological innovation]
\item [Highlight 3: experimental result or impact]
\end{highlights}

% Keywords
\begin{keywords}
computational biology \sep bioinformatics \sep [keyword3] \sep [keyword4] \sep [keyword5]
\end{keywords}

\maketitle

%% ============================================================
%% Main Text
%% ============================================================

\section{Introduction}
\label{sec:intro}

[Introduction text. Provide background, motivation, and state the contributions of this work.]

\section{Related Work}
\label{sec:related}

[Review relevant literature. Organize by topic or methodology.]

\section{Methodology}
\label{sec:method}

\subsection{Problem Formulation}
\label{subsec:problem}

[Define the problem formally. Use mathematical notation as needed, \eg:]
% Let $\mathcal{D} = \{(x_i, y_i)\}_{i=1}^{N}$ denote the training set.

\subsection{Proposed Method}
\label{subsec:proposed}

[Describe the proposed method in detail.]

\section{Experiments}
\label{sec:experiments}

\subsection{Experimental Setup}
\label{subsec:setup}

[Describe hardware, software, hyperparameters, and training configuration.]

\subsection{Datasets}
\label{subsec:datasets}

[Describe datasets used. Include a summary table.]

% Example table:
% \begin{table}[htbp]
% \caption{Dataset statistics}\label{tab:datasets}
% \begin{tabular*}{\tblwidth}{@{}lccc@{}}
% \toprule
% Dataset & Samples & Features & Classes \\
% \midrule
% Dataset A & 1000 & 128 & 2 \\
% Dataset B & 5000 & 256 & 10 \\
% \bottomrule
% \end{tabular*}
% \end{table}

\subsection{Baselines}
\label{subsec:baselines}

[Describe baseline methods for comparison.]

\subsection{Results}
\label{subsec:results}

[Present results with tables and figures. Include statistical significance tests where appropriate.]

% Example results table:
% \begin{table}[htbp]
% \caption{Comparison with baselines}\label{tab:results}
% \begin{tabular*}{\tblwidth}{@{}lcccc@{}}
% \toprule
% Method & Accuracy & Precision & Recall & F1 \\
% \midrule
% Baseline 1 & 0.85 & 0.83 & 0.87 & 0.85 \\
% Baseline 2 & 0.88 & 0.86 & 0.89 & 0.87 \\
% \textbf{Ours} & \textbf{0.92} & \textbf{0.91} & \textbf{0.93} & \textbf{0.92} \\
% \bottomrule
% \end{tabular*}
% \end{table}

\section{Discussion}
\label{sec:discussion}

[Discuss implications, limitations, and future directions.]

\section{Conclusion}
\label{sec:conclusion}

[Summarize contributions and findings.]

% Credit authorship contribution statement
\printcredits

%% ============================================================
%% Back Matter
%% ============================================================

\section*{Declaration of Competing Interest}
The authors declare that they have no known competing financial interests or personal relationships that could have appeared to influence the work reported in this paper.

\section*{Acknowledgements}
[Acknowledgements]

%% ============================================================
%% Bibliography
%% ============================================================

\bibliographystyle{cas-model2-names}
\bibliography{references}

%% ============================================================
%% Appendix (uncomment if needed)
%% ============================================================

% \appendix
% \section{Supplementary Material}
% \label{app:supplementary}

\end{document}
